\documentclass[11pt,a4paper]{article}

% --- Encoding & Sprache ---
\usepackage[utf8]{inputenc}
\usepackage[T1]{fontenc}
\usepackage[ngerman]{babel}
\usepackage{lmodern}

% --- Layout ---
\usepackage[margin=2.5cm]{geometry}
\usepackage{parskip}

% --- Mathe, Grafiken, Tabellen ---
\usepackage{amsmath, amssymb}
\usepackage{graphicx}
\usepackage{booktabs}

% --- Code-Listings ---
\usepackage{xcolor}
\usepackage{listings}
\lstset{
  basicstyle=\ttfamily\small,
  keywordstyle=\color{blue},
  commentstyle=\color{gray},
  stringstyle=\color{teal},
  numbers=left,
  numberstyle=\tiny\color{gray},
  frame=single,
  breaklines=true,
  showstringspaces=false
}

% --- Links (zuletzt laden) ---
\usepackage[hidelinks]{hyperref}

\title{Advanced Software Engineering}
\author{Merlin Bürkle, Tom Schneider}
\date{\today}

\begin{document}

\maketitle
\tableofcontents
\newpage

\section{Einleitung}
Hier beginnt der Text.

\subsection{Beispiel: Formel}
\begin{equation}
  E = mc^2
\end{equation}

\subsection{Beispiel: Tabelle}
\begin{table}[h]
  \centering
  \begin{tabular}{lr}
    \toprule
    Spalte A & Spalte B \\
    \midrule
    Wert 1 & 42 \\
    Wert 2 & 23 \\
    \bottomrule
  \end{tabular}
  \caption{Beispieltabelle}
\end{table}

\subsection{Beispiel: Code}
\begin{lstlisting}[language=Java]
public class Hello {
    public static void main(String[] args) {
        System.out.println("Hello World");
    }
}
\end{lstlisting}

\end{document}
